\documentclass[11pt]{article}

\usepackage[margin=2.5cm]{geometry}
\usepackage{amsmath,amssymb,amsthm,bm}
\usepackage{mathtools}

\title{Correction -- Exercise Sheet: Conditioning \& Martingales (Applied)}
\author{}
\date{}

\begin{document}
\maketitle

\paragraph{Conventions.} Throughout, all equalities involving conditional expectations hold almost surely.

% =====================================================
\section*{Part A -- Fundamental and Applied Exercises}
% =====================================================

\begin{enumerate}

% -------------------------------------------------
\item \textbf{Conditional expectation with a partition}

\begin{enumerate}
\item Since $\mathbb{P}(\omega_i)=1/4$,
\[
\mathbb{E}[X]=\frac14(1+3+0+2)=\frac64=\frac32.
\]

\item Let $\mathcal{G}=\sigma(B_1,B_2)$. Any $\mathcal{G}$-measurable r.v.\ is constant on $B_1$ and on $B_2$.
Thus
\[
\mathbb{E}[X\mid\mathcal{G}]
= c_1\mathbf{1}_{B_1}+c_2\mathbf{1}_{B_2},
\qquad
c_i=\mathbb{E}[X\mid B_i].
\]
Compute
\[
c_1=\frac{\mathbb{E}[X\mathbf{1}_{B_1}]}{\mathbb{P}(B_1)}
=\frac{\frac14(1+3)}{\frac12}=4\cdot \frac14\cdot \frac{1+3}{2}=2,
\qquad
c_2=\frac{\frac14(0+2)}{\frac12}=1.
\]
Hence
\[
\boxed{\ \mathbb{E}[X\mid\sigma(B_1,B_2)] = 2\,\mathbf{1}_{B_1}+1\,\mathbf{1}_{B_2}\ }.
\]

\item Averaging:
\[
\mathbb{E}\big[\mathbb{E}[X\mid\mathcal{G}]\big]
=2\,\mathbb{P}(B_1)+1\,\mathbb{P}(B_2)
=2\cdot \tfrac12+1\cdot\tfrac12=\tfrac32=\mathbb{E}[X].
\]
\end{enumerate}

\bigskip

% -------------------------------------------------
\item \textbf{Conditioning on a discrete random variable}

\begin{enumerate}
\item For $i\in\{0,1\}$,
\[
\mathbb{P}(X=i)=\sum_{j=0}^3 \mathbb{P}(X=i,Y=j)=4\cdot \frac18=\frac12.
\]
For $j\in\{0,1,2,3\}$,
\[
\mathbb{P}(Y=j)=\sum_{i=0}^1 \mathbb{P}(X=i,Y=j)=2\cdot \frac18=\frac14.
\]

\item For any $j$ with $\mathbb{P}(Y=j)>0$,
\[
\mathbb{E}[X\mid Y=j]
=\sum_{i\in\{0,1\}} i\,\mathbb{P}(X=i\mid Y=j)
= \mathbb{P}(X=1\mid Y=j).
\]
But
\[
\mathbb{P}(X=1\mid Y=j)=\frac{\mathbb{P}(X=1,Y=j)}{\mathbb{P}(Y=j)}
=\frac{\frac18}{\frac14}=\frac12.
\]
So $\boxed{\mathbb{E}[X\mid Y]=\frac12}$ (a.s.).

\item Check independence:
\[
\mathbb{P}(X=i,Y=j)=\frac18
=\mathbb{P}(X=i)\mathbb{P}(Y=j)=\frac12\cdot \frac14=\frac18.
\]
Hence $X$ and $Y$ are independent.
\end{enumerate}

\bigskip

% -------------------------------------------------
\item \textbf{Conditional expectation as best predictor}

\begin{enumerate}
\item By definition, $\mathbb{E}[X\mid Y]$ is $\sigma(Y)$-measurable, hence there exists a (Borel) function $h$
such that $\mathbb{E}[X\mid Y]=h(Y)$. If $X\in L^2$, then $\mathbb{E}[X\mid Y]\in L^2$ as well, so $h(Y)\in L^2(Y)$.

\item Let $g(Y)\in L^2(Y)$. Using the defining property with the test function $g(Y)$ (or the tower property),
\[
\mathbb{E}\!\left[(X-\mathbb{E}[X\mid Y])\,g(Y)\right]
=\mathbb{E}\!\left[\ \mathbb{E}\!\left[(X-\mathbb{E}[X\mid Y])\,g(Y)\mid Y\right]\right].
\]
Since $g(Y)$ is $\sigma(Y)$-measurable,
\[
\mathbb{E}\!\left[(X-\mathbb{E}[X\mid Y])\,g(Y)\mid Y\right]
=g(Y)\,\mathbb{E}\!\left[X-\mathbb{E}[X\mid Y]\mid Y\right]
=g(Y)\,\big(\mathbb{E}[X\mid Y]-\mathbb{E}[X\mid Y]\big)=0.
\]
So the expectation is $0$.

\item In the Hilbert space $L^2$, the subspace $L^2(Y)$ is the set of square-integrable functions of $Y$.
The identity in (b) says the error $X-\mathbb{E}[X\mid Y]$ is orthogonal to $L^2(Y)$, i.e.\ $\mathbb{E}[X\mid Y]$
is the orthogonal projection of $X$ onto $L^2(Y)$, hence the best mean-square predictor using $Y$.
\end{enumerate}

\bigskip

% -------------------------------------------------
\item \textbf{Gaussian conditioning (one-dimensional regression)}

Here $(X,Y)$ is centered Gaussian with $\mathrm{Var}(X)=1$, $\mathrm{Var}(Y)=4$, $\mathrm{Cov}(X,Y)=1$.

\begin{enumerate}
\item For jointly Gaussian variables,
\[
\mathbb{E}[Y\mid X]=\mathbb{E}[Y]+\frac{\mathrm{Cov}(Y,X)}{\mathrm{Var}(X)}(X-\mathbb{E}[X])
=\frac{1}{1}\,X=X.
\]

\item The conditional variance formula gives
\[
\mathrm{Var}(Y\mid X)=\mathrm{Var}(Y)-\frac{\mathrm{Cov}(Y,X)^2}{\mathrm{Var}(X)}
=4-\frac{1^2}{1}=3.
\]

\item Regression interpretation: the best linear predictor of $Y$ given $X$ is $Y\approx X$, and the residual
$Y-\mathbb{E}[Y\mid X]$ has variance $3$ (the ``unexplained'' part).
\end{enumerate}

\bigskip

% -----------------------------
\item \textbf{Conditional expectation of an indicator}

\begin{enumerate}
\item Conditioning on an event $B$ is the same as conditioning on the $\sigma$-algebra generated by $B$ restricted to $B$:
\[
\mathbb{E}[\mathbf{1}_A\mid B]
:=\mathbb{E}[\mathbf{1}_A\mid \{B\ \text{occurs}\}]
=\frac{\mathbb{E}[\mathbf{1}_A\mathbf{1}_B]}{\mathbb{P}(B)}
=\frac{\mathbb{P}(A\cap B)}{\mathbb{P}(B)}
=\mathbb{P}(A\mid B).
\]

\item Any $\sigma(B)$-measurable r.v.\ has the form $c_1\mathbf{1}_B+c_0\mathbf{1}_{B^c}$.
Using the characterization on $B$ and $B^c$,
\[
\mathbb{E}[\mathbf{1}_A\mid \sigma(B)]
=\mathbb{E}[\mathbf{1}_A\mid B]\mathbf{1}_B+\mathbb{E}[\mathbf{1}_A\mid B^c]\mathbf{1}_{B^c}
=\mathbb{P}(A\mid B)\mathbf{1}_B+\mathbb{P}(A\mid B^c)\mathbf{1}_{B^c}.
\]

\item Interpreting $\mathbf{1}_A$ as the payoff of a digital claim that pays $1$ if $A$ happens, the conditional
expectation is the ``best forecast'' of that payoff given the information ``whether $B$ happened''.
\end{enumerate}

\bigskip

% -----------------------------
\item \textbf{Conditional expectation and independence}

\begin{enumerate}
\item If $X$ is independent of $Y$, then for any bounded measurable $g$,
\[
\mathbb{E}[Xg(Y)]=\mathbb{E}[X]\mathbb{E}[g(Y)].
\]
Let $Z=\mathbb{E}[X\mid Y]$. By the defining property of conditional expectation,
\[
\mathbb{E}[Zg(Y)]=\mathbb{E}[Xg(Y)]=\mathbb{E}[X]\mathbb{E}[g(Y)].
\]
But also $\mathbb{E}[\mathbb{E}[X]\cdot g(Y)]=\mathbb{E}[X]\mathbb{E}[g(Y)]$, hence $Z$ and the constant $\mathbb{E}[X]$
have the same integrals against all $g(Y)$, which implies $Z=\mathbb{E}[X]$ a.s.

\item Example: let $U\sim U[-1,1]$, set $X=U$ and $Y=U^2$. Then $\mathbb{E}[X\mid Y]=0$ a.s. by symmetry,
so $\mathbb{E}[X\mid Y]=\mathbb{E}[X]=0$, but $X$ and $Y$ are not independent (since $Y$ is a function of $X$).

\item Equality $\mathbb{E}[X\mid Y]=\mathbb{E}[X]$ only says $Y$ carries no \emph{linear/mean} predictive information about $X$.
Independence is stronger: it requires factorization of the whole joint law.
\end{enumerate}

\bigskip

% -----------------------------
\item \textbf{Conditional variance decomposition}

\begin{enumerate}
\item Let $m(Y)=\mathbb{E}[X\mid Y]$. Write
\[
X-\mathbb{E}[X] = \big(X-m(Y)\big) + \big(m(Y)-\mathbb{E}[X]\big).
\]
Square and take expectations:
\begin{align*}
\mathrm{Var}(X)
&=\mathbb{E}\big[(X-\mathbb{E}[X])^2\big]\\
&=\mathbb{E}\big[(X-m(Y))^2\big]
+2\,\mathbb{E}\big[(X-m(Y))(m(Y)-\mathbb{E}[X])\big]
+\mathbb{E}\big[(m(Y)-\mathbb{E}[X])^2\big].
\end{align*}
The cross term is $0$ because $m(Y)-\mathbb{E}[X]$ is $\sigma(Y)$-measurable and
$\mathbb{E}[X-m(Y)\mid Y]=0$:
\[
\mathbb{E}\big[(X-m(Y))(m(Y)-\mathbb{E}[X])\big]
=\mathbb{E}\Big[\ (m(Y)-\mathbb{E}[X])\,\mathbb{E}[X-m(Y)\mid Y]\ \Big]=0.
\]
Also, by definition $\mathrm{Var}(X\mid Y):=\mathbb{E}\big[(X-m(Y))^2\mid Y\big]$, hence
\[
\mathbb{E}\big[(X-m(Y))^2\big]=\mathbb{E}\big[\mathrm{Var}(X\mid Y)\big],
\qquad
\mathbb{E}\big[(m(Y)-\mathbb{E}[X])^2\big]=\mathrm{Var}(m(Y)).
\]
Therefore
\[
\boxed{\ \mathrm{Var}(X)=\mathbb{E}[\mathrm{Var}(X\mid Y)] + \mathrm{Var}(\mathbb{E}[X\mid Y])\ }.
\]

\item $\mathbb{E}[\mathrm{Var}(X\mid Y)]$ is the average remaining uncertainty after observing $Y$ (unexplained part),
while $\mathrm{Var}(\mathbb{E}[X\mid Y])$ is the variability of the predictor itself (explained part).

\item In regression, $m(Y)$ is the fitted value and $X-m(Y)$ is the residual. The identity is the classical
``total = explained + residual'' variance decomposition.
\end{enumerate}

\bigskip

% -----------------------------
\item \textbf{Convergence in probability: basic examples}

\begin{enumerate}
\item Fix $\varepsilon>0$. Then
\[
\mathbb{P}(|X_n|>\varepsilon)=\mathbb{P}\left(\frac{|Z|}{n}>\varepsilon\right)
=\mathbb{P}\left(|Z|>n\varepsilon\right)\xrightarrow[n\to\infty]{}0,
\]
since $\{|Z|>n\varepsilon\}\downarrow \emptyset$ and probabilities are continuous from above.
Hence $X_n\xrightarrow{\mathbb{P}}0$.

\item Since $Z\in L^1$,
\[
\mathbb{E}[|X_n|]=\mathbb{E}\left[\frac{|Z|}{n}\right]=\frac{1}{n}\mathbb{E}[|Z|]\to 0,
\]
so $X_n\xrightarrow{L^1}0$.

\item Yes: almost sure convergence holds without extra assumptions. Indeed, for every $\omega$,
$X_n(\omega)=Z(\omega)/n\to 0$. Thus $X_n\to 0$ pointwise, hence a.s.
\end{enumerate}

\end{enumerate}

\bigskip
\hrule
\bigskip

% =====================================================
\section*{Part B -- More Advanced Applied Exercises}
% =====================================================

\begin{enumerate}
\setcounter{enumi}{8}

% -------------------------------------------------
\item \textbf{Martingale check (two-step random walk)}

Let $X_0=0$, and $X_n=X_{n-1}+\xi_n$, where $(\xi_n)$ are i.i.d. with
$\mathbb{P}(\xi_n=1)=\mathbb{P}(\xi_n=-1)=1/2$.

\begin{enumerate}
\item For the natural filtration $\mathcal{F}_n=\sigma(\xi_1,\dots,\xi_n)$,
\[
\mathbb{E}[X_{n+1}\mid \mathcal{F}_n]=\mathbb{E}[X_n+\xi_{n+1}\mid \mathcal{F}_n]
= X_n + \mathbb{E}[\xi_{n+1}]=X_n,
\]
since $\xi_{n+1}$ is independent of $\mathcal{F}_n$ and has mean $0$. Hence $(X_n)$ is a martingale.

\item For $Y_n=X_n^2-n$, use
\[
X_{n+1}^2=(X_n+\xi_{n+1})^2=X_n^2+2X_n\xi_{n+1}+\xi_{n+1}^2.
\]
Conditioning on $\mathcal{F}_n$ gives
\[
\mathbb{E}[X_{n+1}^2\mid\mathcal{F}_n]=X_n^2+2X_n\mathbb{E}[\xi_{n+1}\mid\mathcal{F}_n]+\mathbb{E}[\xi_{n+1}^2]
=X_n^2+1,
\]
because $\mathbb{E}[\xi_{n+1}]=0$ and $\xi_{n+1}^2=1$. Therefore
\[
\mathbb{E}[Y_{n+1}\mid\mathcal{F}_n]
=\mathbb{E}[X_{n+1}^2-(n+1)\mid\mathcal{F}_n]
=(X_n^2+1)-(n+1)=X_n^2-n=Y_n,
\]
so $(Y_n)$ is also a martingale.
\end{enumerate}

\bigskip

% -------------------------------------------------
\item \textbf{Doob decomposition (submartingale)}

Let $X_n=\sum_{k=1}^n \xi_k$ as above and define $M_n=X_n^2-n$.

\begin{enumerate}
\item Since $(M_n)$ is a martingale (previous exercise), $X_n^2 = M_n + n$ is a decomposition into a martingale plus a predictable increasing process. Moreover, $n$ is increasing and deterministic, hence predictable.

\item To see that $(X_n^2)$ is a submartingale, compute
\[
\mathbb{E}[X_{n+1}^2\mid \mathcal{F}_n]=X_n^2+1 \ge X_n^2.
\]
So $\mathbb{E}[X_{n+1}^2\mid\mathcal{F}_n]\ge X_n^2$, i.e. submartingale.
\end{enumerate}

\bigskip

% -------------------------------------------------
\item \textbf{One-step binomial pricing (mini finance)}

Here $S_0=100$ and $S_1\in\{120,90\}$, with risk-free rate $r=0$.

\begin{enumerate}
\item Under a risk-neutral measure $\mathbb{Q}$, discounted prices are martingales. Since $r=0$, discounting is $1$ and we require
\[
\mathbb{E}^{\mathbb{Q}}[S_1]=S_0.
\]
Let $q=\mathbb{Q}(S_1=120)$. Then $\mathbb{Q}(S_1=90)=1-q$ and
\[
100 = 120q+90(1-q)=90+30q \quad\Rightarrow\quad q=\frac{10}{30}=\frac13.
\]
Thus $\boxed{q=1/3}$.

\item The payoff is $H=(S_1-100)_+$, so $H=20$ in the up state and $H=0$ in the down state. Hence
\[
V_0=\mathbb{E}^{\mathbb{Q}}[H]=20\cdot q+0\cdot(1-q)=20\cdot\frac13=\frac{20}{3}.
\]

\item Replication: find $(\Delta,B)$ such that $\Delta S_1+B=H$ in both states:
\[
120\Delta+B=20,\qquad 90\Delta+B=0.
\]
Subtract: $30\Delta=20$ so $\Delta=2/3$. Then $B=-90\Delta=-60$.
Initial cost is $\Delta S_0+B=\frac23\cdot 100-60=\frac{200}{3}-\frac{180}{3}=\frac{20}{3}$, matching $V_0$.
\end{enumerate}

\bigskip

% -------------------------------------------------
\item \textbf{Optional stopping on a bounded stopping time (simple)}

Let $(M_n)$ be a martingale and $\tau$ a bounded stopping time, i.e. $\tau\le N$ a.s.

\begin{enumerate}
\item A standard (bounded) optional stopping result states:
if $\tau$ is bounded, then
\[
\mathbb{E}[M_\tau]=\mathbb{E}[M_0].
\]
Sketch: write $M_\tau=\sum_{k=0}^N M_k\mathbf{1}_{\{\tau=k\}}$, take expectations and use the tower property plus martingale increments, or apply Doob's optional sampling theorem (bounded case).

\item Apply to the martingale $M_n=X_n$ from the random walk: for any bounded stopping time $\tau$, $\mathbb{E}[X_\tau]=0$.

\item In finance language: a ``fair game'' has zero expected gain even if you stop at a (bounded) random time based on past information.
\end{enumerate}

\end{enumerate}

\end{document}
